%% Generated by tools/gen_error_codes.py from docs/errors/ of the compiler;
%% edit the compiler's pages or tools/error_themes.txt, then rerun it.

\clearpage
\section{Functions, calls and closures}
\label{sec:codes:functions}

\begin{longtable}{@{}l p{0.78\linewidth} r@{}}
\toprule Code & Message & Page\\ \midrule \endhead
\hyperref[sec:codes:e4023]{E4023} & colliding function definitions \errhole{} and \errhole{} & \pageref{sec:codes:e4023}\\
\hyperref[sec:codes:e4042]{E4042} & parameter with a default value must have a referenceable name & \pageref{sec:codes:e4042}\\
\hyperref[sec:codes:e4043]{E4043} & parameter with a default value cannot be a reference & \pageref{sec:codes:e4043}\\
\hyperref[sec:codes:e4073]{E4073} & cannot create a function pointer, entities are not copiable & \pageref{sec:codes:e4073}\\
\hyperref[sec:codes:e4074]{E4074} & cannot create a function pointer, it would enclose the record by reference & \pageref{sec:codes:e4074}\\
\hyperref[sec:codes:e4075]{E4075} & using a deep copy instead of a one level copy is prohibited when enclosing a record & \pageref{sec:codes:e4075}\\
\hyperref[sec:codes:e4076]{E4076} & cannot create a function pointer from the throwing function \errhole{} & \pageref{sec:codes:e4076}\\
\hyperref[sec:codes:e4077]{E4077} & cannot create a function pointer from the unsafe function \errhole{} & \pageref{sec:codes:e4077}\\
\hyperref[sec:codes:e4078]{E4078} & cannot create an option delegate from the non throwing function \errhole{} & \pageref{sec:codes:e4078}\\
\hyperref[sec:codes:e4084]{E4084} & cannot create a function pointer from \errhole{} implicitly & \pageref{sec:codes:e4084}\\
\hyperref[sec:codes:e4085]{E4085} & cannot create a delegate from the method \errhole{} implicitly & \pageref{sec:codes:e4085}\\
\hyperref[sec:codes:e4105]{E4105} & \errhole{} is a method and must be called from an instance & \pageref{sec:codes:e4105}\\
\hyperref[sec:codes:e4107]{E4107} & call operator is not defined for value \errhole{} & \pageref{sec:codes:e4107}\\
\hyperref[sec:codes:e4113]{E4113} & main function takes at most one argument & \pageref{sec:codes:e4113}\\
\hyperref[sec:codes:e4114]{E4114} & main function cannot be inline & \pageref{sec:codes:e4114}\\
\hyperref[sec:codes:e4134]{E4134} & named parameter \errhole{} is set multiple times & \pageref{sec:codes:e4134}\\
\hyperref[sec:codes:e4171]{E4171} & no parameter is named \errhole{} & \pageref{sec:codes:e4171}\\
\hyperref[sec:codes:e4212]{E4212} & \errhole{} parameters \errhole{} expected, but \errhole{} \errhole{} provided & \pageref{sec:codes:e4212}\\
\hyperref[sec:codes:e4226]{E4226} & the call operator is not defined for \errhole{} and \errhole{} & \pageref{sec:codes:e4226}\\
\hyperref[sec:codes:e4246]{E4246} & the creation of a delegate from a method has no effect & \pageref{sec:codes:e4246}\\
\hyperref[sec:codes:e4250]{E4250} & copy the lambda closure has no effect, no values are enclosed & \pageref{sec:codes:e4250}\\
\hyperref[sec:codes:e4280]{E4280} & \errhole{} called with \errhole{} works with multiple candidates & \pageref{sec:codes:e4280}\\
\hyperref[sec:codes:e4283]{E4283} & the parameter \errhole{} has a default value, it can only be passed by name & \pageref{sec:codes:e4283}\\
\hyperref[sec:codes:e4284]{E4284} & the parameter \errhole{} can only be passed by name & \pageref{sec:codes:e4284}\\
\hyperref[sec:codes:e4285]{E4285} & parameter with a default value cannot be marked \errhole{}, it can already only be passed by name & \pageref{sec:codes:e4285}\\
\hyperref[sec:codes:e4298]{E4298} & the method \errhole{} is not a value, it can only be called & \pageref{sec:codes:e4298}\\
\bottomrule
\end{longtable}

\errorcode{E4023}{colliding function definitions \errhole{} and \errhole{}}
\label{sec:codes:e4023}

Raised during \textbf{validation}, from \texttt{ValidateErrorMessage::\allowbreak{}COLLIDING\_\allowbreak{}FUNCTION\_\allowbreak{}DEFINITION} (\textit{src/\allowbreak{}ymirc/\allowbreak{}semantic/\allowbreak{}validator/\allowbreak{}errors.yr}).

\subsubsection*{What it means}

Two functions have the same name and signatures that overload resolution cannot tell apart, so a call naming them would be ambiguous.

\subsubsection*{Example}

\textit{test\_\allowbreak{}resources/\allowbreak{}function/\allowbreak{}test15.yr}:

%% check: skip
\begin{lstlisting}[style=coloredverbatimError]
fn foo () {}
fn foo () {}

fn bar (i : i32) {i;}
fn bar () {}


fn baz (i : i32, j : i32) { i; j; }
fn baz (i : i32, ref j : i32) {
    i; j;
}
\end{lstlisting}

\begin{lstlisting}[style=bashVerb, breaklines=true, escapechar=~]
Error[E4023] : colliding function definitions test15::foo ()-> void and test15::foo ()-> void
 --> test_resources/function/test15.yr:(1,4)
 1  ~\lstbox{\textSFxi{}}~ fn foo () {}
    ~\lstbox{\textSFv{}}~    ^^^
   ...
 --> test_resources/function/test15.yr:(2,4)
 2  ~\lstbox{\textSFxi{}}~ fn foo () {}
    ~\lstbox{\textSFv{}}~    ^^^
    ~\lstbox{\textSFii{}}~~\lstbox{\textSFx{}}~~\lstbox{\textSFx{}}~~\lstbox{\textSFx{}}~~\lstbox{\textSFx{}}~~\lstbox{\textSFx{}}~~\lstbox{\textSFx{}}~~\lstbox{\textSFx{}}~
\end{lstlisting}

\subsubsection*{How to fix it}

Rename one of them, or change the parameters so the two differ in a way a call can select on.

\errorcode{E4042}{parameter with a default value must have a referenceable name}
\label{sec:codes:e4042}

Raised during \textbf{validation}, from \texttt{ValidateErrorMessage::\allowbreak{}DEFAULT\_\allowbreak{}VAR\_\allowbreak{}NO\_\allowbreak{}NAME} (\textit{src/\allowbreak{}ymirc/\allowbreak{}semantic/\allowbreak{}validator/\allowbreak{}errors.yr}).

\subsubsection*{What it means}

A parameter has a default value but no name a caller could use. Default values are only reachable by naming the parameter at the call site, so an unnamed parameter could never take its default.

\subsubsection*{Example}

\textit{test\_\allowbreak{}resources/\allowbreak{}diagnostics/\allowbreak{}test27.yr}:

%% check: skip
\begin{lstlisting}[style=coloredverbatimError]
fn foo (_: i32 = 0)-> i32 { 1 }

fn main () { foo (); }
\end{lstlisting}

\begin{lstlisting}[style=bashVerb, breaklines=true, escapechar=~]
Error[E4042] : parameter with a default value must have a referenceable name
 --> test_resources/diagnostics/test27.yr:(1,9)
 1  ~\lstbox{\textSFxi{}}~ fn foo (_: i32 = 0)-> i32 { 1 }
    ~\lstbox{\textSFv{}}~         ^        -
    ~\lstbox{\textSFxi{}}~
    = when validating test27::foo -- (test_resources/diagnostics/test27.yr:(1,4))
    ~\lstbox{\textSFii{}}~~\lstbox{\textSFx{}}~~\lstbox{\textSFx{}}~~\lstbox{\textSFx{}}~~\lstbox{\textSFx{}}~~\lstbox{\textSFx{}}~~\lstbox{\textSFx{}}~~\lstbox{\textSFx{}}~
Error[E4104] : symbol test27::foo has errors
 --> test_resources/diagnostics/test27.yr:(3,14)
 1  ~\lstbox{\textSFxi{}}~ fn foo (_: i32 = 0)-> i32 { 1 }
    ~\lstbox{\textSFv{}}~         -        -  error: parameter with a default value must have a referenceable name
 2  ~\lstbox{\textSFxi{}}~
 3  ~\lstbox{\textSFxi{}}~ fn main () { foo (); }
    ~\lstbox{\textSFv{}}~              ^^^
    ~\lstbox{\textSFxi{}}~
    = when validating test27::foo -- (test_resources/diagnostics/test27.yr:(1,4))
    = when validating test27::main -- (test_resources/diagnostics/test27.yr:(3,4))
    ~\lstbox{\textSFii{}}~~\lstbox{\textSFx{}}~~\lstbox{\textSFx{}}~~\lstbox{\textSFx{}}~~\lstbox{\textSFx{}}~~\lstbox{\textSFx{}}~~\lstbox{\textSFx{}}~~\lstbox{\textSFx{}}~
\end{lstlisting}

\subsubsection*{How to fix it}

Give the parameter a name, or remove the default value.

\errorcode{E4043}{parameter with a default value cannot be a reference}
\label{sec:codes:e4043}

Raised during \textbf{validation}, from \texttt{ValidateErrorMessage::\allowbreak{}DEFAULT\_\allowbreak{}VAR\_\allowbreak{}REF} (\textit{src/\allowbreak{}ymirc/\allowbreak{}semantic/\allowbreak{}validator/\allowbreak{}errors.yr}).

\subsubsection*{What it means}

A parameter is declared \tokennolst{ref} and given a default value. A reference parameter borrows the caller's storage; a default value has no storage in the caller to borrow.

\subsubsection*{Example}

\textit{test\_\allowbreak{}resources/\allowbreak{}diagnostics/\allowbreak{}test28.yr}:

%% check: skip
\begin{lstlisting}[style=coloredverbatimError]
fn foo (ref x: i32 = 0)-> i32 { x }

fn main () { foo (); }
\end{lstlisting}

\begin{lstlisting}[style=bashVerb, breaklines=true, escapechar=~]
Error[E4043] : parameter with a default value cannot be a reference
 --> test_resources/diagnostics/test28.yr:(1,9)
 1  ~\lstbox{\textSFxi{}}~ fn foo (ref x: i32 = 0)-> i32 { x }
    ~\lstbox{\textSFv{}}~         ^^^ -        -
    ~\lstbox{\textSFxi{}}~
    = when validating test28::foo -- (test_resources/diagnostics/test28.yr:(1,4))
    ~\lstbox{\textSFii{}}~~\lstbox{\textSFx{}}~~\lstbox{\textSFx{}}~~\lstbox{\textSFx{}}~~\lstbox{\textSFx{}}~~\lstbox{\textSFx{}}~~\lstbox{\textSFx{}}~~\lstbox{\textSFx{}}~
Error[E4104] : symbol test28::foo has errors
 --> test_resources/diagnostics/test28.yr:(3,14)
 1  ~\lstbox{\textSFxi{}}~ fn foo (ref x: i32 = 0)-> i32 { x }
    ~\lstbox{\textSFv{}}~         ---          -  error: parameter with a default value cannot be a reference
 2  ~\lstbox{\textSFxi{}}~
 3  ~\lstbox{\textSFxi{}}~ fn main () { foo (); }
    ~\lstbox{\textSFv{}}~              ^^^
    ~\lstbox{\textSFxi{}}~
    = when validating test28::foo -- (test_resources/diagnostics/test28.yr:(1,4))
    = when validating test28::main -- (test_resources/diagnostics/test28.yr:(3,4))
    ~\lstbox{\textSFii{}}~~\lstbox{\textSFx{}}~~\lstbox{\textSFx{}}~~\lstbox{\textSFx{}}~~\lstbox{\textSFx{}}~~\lstbox{\textSFx{}}~~\lstbox{\textSFx{}}~~\lstbox{\textSFx{}}~
\end{lstlisting}

\subsubsection*{How to fix it}

Drop either the \tokennolst{ref} or the default value.

\errorcode{E4073}{cannot create a function pointer, entities are not copiable}
\label{sec:codes:e4073}

Raised during \textbf{validation}, from \texttt{ValidateErrorMessage::\allowbreak{}FUNCTION\_\allowbreak{}TO\_\allowbreak{}FPTR\_\allowbreak{}ENTITY} (\textit{src/\allowbreak{}ymirc/\allowbreak{}semantic/\allowbreak{}validator/\allowbreak{}errors.yr}).

\subsubsection*{What it means}

A function pointer would have to capture an entity, and entities are not copiable, so there is no way to store one in the closure environment.

\subsubsection*{Example}

\textit{test\_\allowbreak{}resources/\allowbreak{}diagnostics/\allowbreak{}test103.yr}:

%% check: skip
\begin{lstlisting}[style=coloredverbatimError]
// a function pointer made from the method of an entity
entity Counter {
    pub self () {}
    pub fn get (self)-> i32 { 1 }
}

fn main () {
    let c = Counter ();
    let _ = &c.get;
}
\end{lstlisting}

\begin{lstlisting}[style=bashVerb, breaklines=true, escapechar=~]
Error[E4073] : cannot create a function pointer, entities are not copiable
 --> test_resources/diagnostics/test103.yr:(9,13)
 9  ~\lstbox{\textSFxi{}}~     let _ = &c.get;
    ~\lstbox{\textSFv{}}~             ^
    ~\lstbox{\textSFxi{}}~
    = when validating test103::main -- (test_resources/diagnostics/test103.yr:(7,4))
    ~\lstbox{\textSFii{}}~~\lstbox{\textSFx{}}~~\lstbox{\textSFx{}}~~\lstbox{\textSFx{}}~~\lstbox{\textSFx{}}~~\lstbox{\textSFx{}}~~\lstbox{\textSFx{}}~~\lstbox{\textSFx{}}~
\end{lstlisting}

\subsubsection*{How to fix it}

Pass the entity as a parameter of the function rather than capturing it.

\errorcode{E4074}{cannot create a function pointer, it would enclose the record by reference}
\label{sec:codes:e4074}

Raised during \textbf{validation}, from \texttt{ValidateErrorMessage::\allowbreak{}FUNCTION\_\allowbreak{}TO\_\allowbreak{}FPTR\_\allowbreak{}RECORD} (\textit{src/\allowbreak{}ymirc/\allowbreak{}semantic/\allowbreak{}validator/\allowbreak{}errors.yr}).

\subsubsection*{What it means}

A function pointer would capture a record by reference, which would outlive the record.

\subsubsection*{Example}

\textit{test\_\allowbreak{}resources/\allowbreak{}diagnostics/\allowbreak{}test104.yr}:

%% check: skip
\begin{lstlisting}[style=coloredverbatimError]
// a function pointer made from the method of a record, outside a copy
record Counter {
    pub self () {}
    pub fn get (self)-> i32 { 1 }
}

fn main () {
    let c = Counter ();
    let _ = &c.get;
}
\end{lstlisting}

\begin{lstlisting}[style=bashVerb, breaklines=true, escapechar=~]
Error[E4074] : cannot create a function pointer, it would enclose the record by reference
 --> test_resources/diagnostics/test104.yr:(9,13)
 9  ~\lstbox{\textSFxi{}}~     let _ = &c.get;
    ~\lstbox{\textSFv{}}~             ^  note: did you mean to enclose it within a copy?
    ~\lstbox{\textSFxi{}}~
    = when validating test104::main -- (test_resources/diagnostics/test104.yr:(7,4))
    ~\lstbox{\textSFii{}}~~\lstbox{\textSFx{}}~~\lstbox{\textSFx{}}~~\lstbox{\textSFx{}}~~\lstbox{\textSFx{}}~~\lstbox{\textSFx{}}~~\lstbox{\textSFx{}}~~\lstbox{\textSFx{}}~
\end{lstlisting}

\subsubsection*{How to fix it}

Capture a copy of the record, or pass it as a parameter.

\errorcode{E4075}{using a deep copy instead of a one level copy is prohibited when enclosing a record}
\label{sec:codes:e4075}

Raised during \textbf{validation}, from \texttt{ValidateErrorMessage::\allowbreak{}FUNCTION\_\allowbreak{}TO\_\allowbreak{}FPTR\_\allowbreak{}RECORD\_\allowbreak{}DCOPY} (\textit{src/\allowbreak{}ymirc/\allowbreak{}semantic/\allowbreak{}validator/\allowbreak{}errors.yr}).

\subsubsection*{What it means}

A record was captured with \tokennolst{dcopy} where the one level \tokennolst{copy} is required, cf. \hyperref[sec:codes:retired]{E4033}.

\subsubsection*{Example}

\textit{test\_\allowbreak{}resources/\allowbreak{}diagnostics/\allowbreak{}test105.yr}:

%% check: skip
\begin{lstlisting}[style=coloredverbatimError]
// a function pointer made from the method of a record, inside a deep copy
record Counter {
    pub self () {}
    pub fn get (self)-> i32 { 1 }
}

fn main () {
    let c = Counter ();
    let _ = dcopy &c.get;
}
\end{lstlisting}

\begin{lstlisting}[style=bashVerb, breaklines=true, escapechar=~]
Error[E4075] : using a deep copy instead of a one level copy is prohibited when enclosing a record
 --> test_resources/diagnostics/test105.yr:(9,19)
 9  ~\lstbox{\textSFxi{}}~     let _ = dcopy &c.get;
    ~\lstbox{\textSFv{}}~                   ^
    ~\lstbox{\textSFxi{}}~
    = help: replace dcopy by the simple copy operator
    ~\lstbox{\textSFxi{}}~
 9  -     let _ = dcopy &c.get;
 9  +     let _ = copy &c.get;
    ~\lstbox{\textSFxi{}}~
    = when validating test105::main -- (test_resources/diagnostics/test105.yr:(7,4))
    ~\lstbox{\textSFii{}}~~\lstbox{\textSFx{}}~~\lstbox{\textSFx{}}~~\lstbox{\textSFx{}}~~\lstbox{\textSFx{}}~~\lstbox{\textSFx{}}~~\lstbox{\textSFx{}}~~\lstbox{\textSFx{}}~
\end{lstlisting}

\subsubsection*{How to fix it}

Use \tokennolst{copy}.

\errorcode{E4076}{cannot create a function pointer from the throwing function \errhole{}}
\label{sec:codes:e4076}

Raised during \textbf{validation}, from \texttt{ValidateErrorMessage::\allowbreak{}FUNCTION\_\allowbreak{}TO\_\allowbreak{}FPTR\_\allowbreak{}THROWING} (\textit{src/\allowbreak{}ymirc/\allowbreak{}semantic/\allowbreak{}validator/\allowbreak{}errors.yr}).

\subsubsection*{What it means}

A function pointer was created from a function that can throw. A function pointer has no place to declare what it throws, so a caller could not know it has to handle an exception.

\subsubsection*{Example}

\textit{test\_\allowbreak{}resources/\allowbreak{}function/\allowbreak{}test11.yr}:

%% check: skip
\begin{lstlisting}[style=coloredverbatimError]
fn foo ()
    throws AssertError
{
    assert (false);
}

fn main () {
    let a = &foo;
    a ();
}
\end{lstlisting}

\begin{lstlisting}[style=bashVerb, breaklines=true, escapechar=~]
Error[E4076] : cannot create a function pointer from the throwing function test11::foo ()-> void
 --> test_resources/function/test11.yr:(8,13)
 8  ~\lstbox{\textSFxi{}}~     let a = &foo;
    ~\lstbox{\textSFv{}}~             ^
    ~\lstbox{\textSFxi{}}~ Note : throwing core::exception::assertion::AssertError
    ~\lstbox{\textSFxi{}}~  --> test_resources/function/test11.yr:(2,12)
    ~\lstbox{\textSFxi{}}~  2  ~\lstbox{\textSFxi{}}~     throws AssertError
    ~\lstbox{\textSFxi{}}~     ~\lstbox{\textSFv{}}~            ^^^^^^^^^^^
    ~\lstbox{\textSFxi{}}~     ~\lstbox{\textSFii{}}~~\lstbox{\textSFx{}}~~\lstbox{\textSFx{}}~~\lstbox{\textSFx{}}~~\lstbox{\textSFx{}}~~\lstbox{\textSFx{}}~~\lstbox{\textSFx{}}~~\lstbox{\textSFx{}}~
    ~\lstbox{\textSFii{}}~~\lstbox{\textSFx{}}~~\lstbox{\textSFx{}}~~\lstbox{\textSFx{}}~~\lstbox{\textSFx{}}~~\lstbox{\textSFx{}}~~\lstbox{\textSFvii{}}~~\lstbox{\textSFx{}}~
\end{lstlisting}

\subsubsection*{How to fix it}

Wrap the call in a lambda that catches the exception, or use a delegate that carries the throw information.

\errorcode{E4077}{cannot create a function pointer from the unsafe function \errhole{}}
\label{sec:codes:e4077}

Raised during \textbf{validation}, from \texttt{ValidateErrorMessage::\allowbreak{}FUNCTION\_\allowbreak{}TO\_\allowbreak{}FPTR\_\allowbreak{}UNSAFE} (\textit{src/\allowbreak{}ymirc/\allowbreak{}semantic/\allowbreak{}validator/\allowbreak{}errors.yr}).

\subsubsection*{What it means}

A function pointer was created from an unsafe function. Calling through the pointer would leave no trace that the call is unsafe.

\subsubsection*{Example}

\textit{test\_\allowbreak{}resources/\allowbreak{}function/\allowbreak{}test13.yr}:

%% check: skip
\begin{lstlisting}[style=coloredverbatimError]
@unsafe
fn foo ();

fn main () {
    let a = &foo;
    a ();
}
\end{lstlisting}

\begin{lstlisting}[style=bashVerb, breaklines=true, escapechar=~]
Error[E4077] : cannot create a function pointer from the unsafe function test13::foo ()-> void
 --> test_resources/function/test13.yr:(5,13)
 5  ~\lstbox{\textSFxi{}}~     let a = &foo;
    ~\lstbox{\textSFv{}}~             ^
    ~\lstbox{\textSFii{}}~~\lstbox{\textSFx{}}~~\lstbox{\textSFx{}}~~\lstbox{\textSFx{}}~~\lstbox{\textSFx{}}~~\lstbox{\textSFx{}}~~\lstbox{\textSFx{}}~~\lstbox{\textSFx{}}~
\end{lstlisting}

\subsubsection*{How to fix it}

Wrap the call in a safe function that enters an \tokennolst{unsafe} block itself, and take a pointer to that.

\errorcode{E4078}{cannot create an option delegate from the non throwing function \errhole{}}
\label{sec:codes:e4078}

Raised during \textbf{validation}, from \texttt{ValidateErrorMessage::\allowbreak{}FUNCTION\_\allowbreak{}TO\_\allowbreak{}OPTION\_\allowbreak{}NO\_\allowbreak{}THROW} (\textit{src/\allowbreak{}ymirc/\allowbreak{}semantic/\allowbreak{}validator/\allowbreak{}errors.yr}).

\subsubsection*{What it means}

The \tokennolst{\&?} operator, which wraps a throwing function into one returning an option, was applied to a function that does not throw. There is no error for the option to carry.

\subsubsection*{Example}

\textit{test\_\allowbreak{}resources/\allowbreak{}function/\allowbreak{}test19.yr}:

%% check: skip
\begin{lstlisting}[style=coloredverbatimError]
fn foo (i : i32)-> i32 {
    i + 1
}

fn main () {
    let _f_ = &?foo; // nothing to catch, '&' is what creates a pointer on 'foo'
}
\end{lstlisting}

\begin{lstlisting}[style=bashVerb, breaklines=true, escapechar=~]
Error[E4078] : cannot create an option delegate from the non throwing function test19::foo (i: i32)-> i32
 --> test_resources/function/test19.yr:(6,15)
 6  ~\lstbox{\textSFxi{}}~     let _f_ = &?foo; // nothing to catch, '&' is what creates a pointer on 'foo'
    ~\lstbox{\textSFv{}}~               ^^
    ~\lstbox{\textSFxi{}}~
    = help: replace the operator &? by the simple & operator
    ~\lstbox{\textSFxi{}}~
 6  -     let _f_ = &?foo; // nothing to catch, '&' is what creates a pointer on 'foo'
 6  +     let _f_ = &foo; // nothing to catch, '&' is what creates a pointer on 'foo'
    ~\lstbox{\textSFii{}}~~\lstbox{\textSFx{}}~~\lstbox{\textSFx{}}~~\lstbox{\textSFx{}}~~\lstbox{\textSFx{}}~~\lstbox{\textSFx{}}~~\lstbox{\textSFx{}}~~\lstbox{\textSFx{}}~
\end{lstlisting}

\subsubsection*{How to fix it}

Use a plain \tokennolst{\&} to take the address of a non-throwing function.

\errorcode{E4084}{cannot create a function pointer from \errhole{} implicitly}
\label{sec:codes:e4084}

Raised during \textbf{validation}, from \texttt{ValidateErrorMessage::\allowbreak{}IMPLICIT\_\allowbreak{}ADDRESS\_\allowbreak{}FUNCTION} (\textit{src/\allowbreak{}ymirc/\allowbreak{}semantic/\allowbreak{}validator/\allowbreak{}errors.yr}).

\subsubsection*{What it means}

A function was used where a function pointer is expected, without asking for its address. Ymir does not take the address of a function implicitly, so a function name used as a value is a mistake rather than a conversion.

\subsubsection*{Example}

\textit{test\_\allowbreak{}resources/\allowbreak{}diagnostics/\allowbreak{}test39.yr}:

%% check: skip
\begin{lstlisting}[style=coloredverbatimError]
fn foo ()-> i32 { 0 }

fn main () {
    let _: fn ()-> i32 = foo;
}
\end{lstlisting}

\begin{lstlisting}[style=bashVerb, breaklines=true, escapechar=~]
Error[E4084] : cannot create a function pointer from test39::foo ()-> i32 implicitly
 --> test_resources/diagnostics/test39.yr:(1,4)
 1  ~\lstbox{\textSFxi{}}~ fn foo ()-> i32 { 0 }
    ~\lstbox{\textSFv{}}~    ^^^
   ...
 --> test_resources/diagnostics/test39.yr:(4,26)
 4  ~\lstbox{\textSFxi{}}~     let _: fn ()-> i32 = foo;
    ~\lstbox{\textSFv{}}~                          ---  note: did you forget ()?
    ~\lstbox{\textSFxi{}}~
    = help: or add the missing &
    ~\lstbox{\textSFxi{}}~
 4  -     let _: fn ()-> i32 = foo;
 4  +     let _: fn ()-> i32 = &foo;
    ~\lstbox{\textSFii{}}~~\lstbox{\textSFx{}}~~\lstbox{\textSFx{}}~~\lstbox{\textSFx{}}~~\lstbox{\textSFx{}}~~\lstbox{\textSFx{}}~~\lstbox{\textSFx{}}~~\lstbox{\textSFx{}}~
\end{lstlisting}

\subsubsection*{How to fix it}

Write \tokennolst{\&} before the function name.

\errorcode{E4085}{cannot create a delegate from the method \errhole{} implicitly}
\label{sec:codes:e4085}

Raised during \textbf{validation}, from \texttt{ValidateErrorMessage::\allowbreak{}IMPLICIT\_\allowbreak{}ADDRESS\_\allowbreak{}METHOD} (\textit{src/\allowbreak{}ymirc/\allowbreak{}semantic/\allowbreak{}validator/\allowbreak{}errors.yr}).

\subsubsection*{What it means}

A method was used where a delegate is expected, without asking for its address.

\subsubsection*{Example}

\textit{test\_\allowbreak{}resources/\allowbreak{}diagnostics/\allowbreak{}test40.yr}:

%% check: skip
\begin{lstlisting}[style=coloredverbatimError]
class Foo {
    pub self () {}
    pub fn bar (self)-> i32 { 0 }
}

fn main () {
    let f = copy Foo ();
    let _: dg ()-> i32 = f.bar;
}
\end{lstlisting}

\begin{lstlisting}[style=bashVerb, breaklines=true, escapechar=~]
Error[E4085] : cannot create a delegate from the method f.test40::Foo::bar ()-> i32 implicitly
 --> test_resources/diagnostics/test40.yr:(8,27)
 8  ~\lstbox{\textSFxi{}}~     let _: dg ()-> i32 = f.bar;
    ~\lstbox{\textSFv{}}~                           ^  note: did you forget ()?
    ~\lstbox{\textSFxi{}}~
    = help: or add the missing &
    ~\lstbox{\textSFxi{}}~
 8  -     let _: dg ()-> i32 = f.bar;
 8  +     let _: dg ()-> i32 = &f.bar;
    ~\lstbox{\textSFii{}}~~\lstbox{\textSFx{}}~~\lstbox{\textSFx{}}~~\lstbox{\textSFx{}}~~\lstbox{\textSFx{}}~~\lstbox{\textSFx{}}~~\lstbox{\textSFx{}}~~\lstbox{\textSFx{}}~
\end{lstlisting}

\subsubsection*{How to fix it}

Write \tokennolst{\&} before the method access.

\errorcode{E4105}{\errhole{} is a method and must be called from an instance}
\label{sec:codes:e4105}

Raised during \textbf{validation}, from \texttt{ValidateErrorMessage::\allowbreak{}IS\_\allowbreak{}A\_\allowbreak{}METHOD} (\textit{src/\allowbreak{}ymirc/\allowbreak{}semantic/\allowbreak{}validator/\allowbreak{}errors.yr}).

\subsubsection*{What it means}

A method was used as if it were a free function. A method needs the instance it runs on.

\subsubsection*{Example}

\textit{test\_\allowbreak{}resources/\allowbreak{}lit\_\allowbreak{}class/\allowbreak{}errors/\allowbreak{}test1.yr}:

%% check: skip
\begin{lstlisting}[style=coloredverbatimError]
class A {
    pub self () {}

    pub fn bar (self) {}
    pub fn foo (self) {
        bar ();
    }
}
\end{lstlisting}

\begin{lstlisting}[style=bashVerb, breaklines=true, escapechar=~]
Error[E4105] : test1::A::bar is a method and must be called from an instance
 --> test_resources/lit_class/errors/test1.yr:(8,9)
 6  ~\lstbox{\textSFxi{}}~     pub fn bar (self) {}
    ~\lstbox{\textSFv{}}~            ---  note: defined here
 7  ~\lstbox{\textSFxi{}}~     pub fn foo (self) {
 8  ~\lstbox{\textSFxi{}}~         bar ();
    ~\lstbox{\textSFv{}}~         ^^^
    ~\lstbox{\textSFxi{}}~
    = when validating test1::A::foo -- (test_resources/lit_class/errors/test1.yr:(7,12))
    = when validating -- (test_resources/lit_class/errors/test1.yr:(3,7))
    ~\lstbox{\textSFii{}}~~\lstbox{\textSFx{}}~~\lstbox{\textSFx{}}~~\lstbox{\textSFx{}}~~\lstbox{\textSFx{}}~~\lstbox{\textSFx{}}~~\lstbox{\textSFx{}}~~\lstbox{\textSFx{}}~
\end{lstlisting}

\subsubsection*{How to fix it}

Call it on an instance (\tokennolst{x.m ()}), or take a delegate with \tokennolst{\&x.m} if the call is to happen later.

\errorcode{E4107}{call operator is not defined for value \errhole{}}
\label{sec:codes:e4107}

Raised during \textbf{validation}, from \texttt{ValidateErrorMessage::\allowbreak{}IS\_\allowbreak{}NOT\_\allowbreak{}CALLABLE} (\textit{src/\allowbreak{}ymirc/\allowbreak{}semantic/\allowbreak{}validator/\allowbreak{}errors.yr}).

\subsubsection*{What it means}

A value was called and its type defines no call operator.

\subsubsection*{Example}

\textit{test\_\allowbreak{}resources/\allowbreak{}diagnostics/\allowbreak{}test48.yr}:

%% check: skip
\begin{lstlisting}[style=coloredverbatimError]
fn main () {
    let x = 1;
    x ();
}
\end{lstlisting}

\begin{lstlisting}[style=bashVerb, breaklines=true, escapechar=~]
Error[E4107] : call operator is not defined for value 1
 --> test_resources/diagnostics/test48.yr:(2,13)
 2  ~\lstbox{\textSFxi{}}~     let x = 1;
    ~\lstbox{\textSFv{}}~             ^
    ~\lstbox{\textSFii{}}~~\lstbox{\textSFx{}}~~\lstbox{\textSFx{}}~~\lstbox{\textSFx{}}~~\lstbox{\textSFx{}}~~\lstbox{\textSFx{}}~~\lstbox{\textSFx{}}~~\lstbox{\textSFx{}}~
\end{lstlisting}

\subsubsection*{How to fix it}

Check that the value is the function, the function pointer or the delegate intended. A class can be made callable by defining the call operator on it.

\errorcode{E4113}{main function takes at most one argument}
\label{sec:codes:e4113}

Raised during \textbf{validation}, from \texttt{ValidateErrorMessage::\allowbreak{}MAIN\_\allowbreak{}FUNCTION\_\allowbreak{}ONE\_\allowbreak{}ARG} (\textit{src/\allowbreak{}ymirc/\allowbreak{}semantic/\allowbreak{}validator/\allowbreak{}errors.yr}).

\subsubsection*{What it means}

The \tokennolst{main} function declares more than one parameter. It takes either nothing or the command line arguments.

\subsubsection*{Example}

\textit{test\_\allowbreak{}resources/\allowbreak{}diagnostics/\allowbreak{}test141.yr}:

%% check: skip
\begin{lstlisting}[style=coloredverbatimError]
// a main function taking a second parameter
fn main (args : [[c8]], count : i32) {
    let _ = (args, count);
}
\end{lstlisting}

\begin{lstlisting}[style=bashVerb, breaklines=true, escapechar=~]
Error[E4113] : main function takes at most one argument
 --> test_resources/diagnostics/test141.yr:(2,25)
 2  ~\lstbox{\textSFxi{}}~ fn main (args : [[c8]], count : i32) {
    ~\lstbox{\textSFv{}}~                         ^^^^^
    ~\lstbox{\textSFxi{}}~
    = when validating test141::main -- (test_resources/diagnostics/test141.yr:(2,4))
    ~\lstbox{\textSFii{}}~~\lstbox{\textSFx{}}~~\lstbox{\textSFx{}}~~\lstbox{\textSFx{}}~~\lstbox{\textSFx{}}~~\lstbox{\textSFx{}}~~\lstbox{\textSFx{}}~~\lstbox{\textSFx{}}~
\end{lstlisting}

\subsubsection*{How to fix it}

Reduce it to at most one parameter.

\errorcode{E4114}{main function cannot be inline}
\label{sec:codes:e4114}

Raised during \textbf{validation}, from \texttt{ValidateErrorMessage::\allowbreak{}MAIN\_\allowbreak{}INLINE} (\textit{src/\allowbreak{}ymirc/\allowbreak{}semantic/\allowbreak{}validator/\allowbreak{}errors.yr}).

\subsubsection*{What it means}

The \tokennolst{main} function is marked \tokennolst{inline}. It is the entry point the runtime calls, so it has to exist as a symbol.

\subsubsection*{Example}

\textit{test\_\allowbreak{}resources/\allowbreak{}diagnostics/\allowbreak{}test49.yr}:

%% check: skip
\begin{lstlisting}[style=coloredverbatimError]
@inline
fn main () {}
\end{lstlisting}

\begin{lstlisting}[style=bashVerb, breaklines=true, escapechar=~]
Error[E4114] : main function cannot be inline
 --> test_resources/diagnostics/test49.yr:(1,2)
 1  ~\lstbox{\textSFxi{}}~ @inline
    ~\lstbox{\textSFv{}}~  ^^^^^^
    ~\lstbox{\textSFii{}}~~\lstbox{\textSFx{}}~~\lstbox{\textSFx{}}~~\lstbox{\textSFx{}}~~\lstbox{\textSFx{}}~~\lstbox{\textSFx{}}~~\lstbox{\textSFx{}}~~\lstbox{\textSFx{}}~
\end{lstlisting}

\subsubsection*{How to fix it}

Remove \tokennolst{inline}.

\errorcode{E4134}{named parameter \errhole{} is set multiple times}
\label{sec:codes:e4134}

Raised during \textbf{validation}, from \texttt{ValidateErrorMessage::\allowbreak{}MULTIPLE\_\allowbreak{}NAMED\_\allowbreak{}PARAM} (\textit{src/\allowbreak{}ymirc/\allowbreak{}semantic/\allowbreak{}validator/\allowbreak{}errors.yr}).

\subsubsection*{What it means}

The same parameter was given twice by name at a call site.

\subsubsection*{Example}

\textit{test\_\allowbreak{}resources/\allowbreak{}diagnostics/\allowbreak{}test53.yr}:

%% check: skip
\begin{lstlisting}[style=coloredverbatimError]
fn foo (a: i32 = 0, b: i32 = 0)-> i32 { a + b }

fn main () { foo (a-> 1, a-> 2); }
\end{lstlisting}

\begin{lstlisting}[style=bashVerb, breaklines=true, escapechar=~]
Error[E4134] : named parameter a is set multiple times
 --> test_resources/diagnostics/test53.yr:(3,23)
 3  ~\lstbox{\textSFxi{}}~ fn main () { foo (a-> 1, a-> 2); }
    ~\lstbox{\textSFv{}}~                       ^   ^^
    ~\lstbox{\textSFii{}}~~\lstbox{\textSFx{}}~~\lstbox{\textSFx{}}~~\lstbox{\textSFx{}}~~\lstbox{\textSFx{}}~~\lstbox{\textSFx{}}~~\lstbox{\textSFx{}}~~\lstbox{\textSFx{}}~
\end{lstlisting}

\subsubsection*{How to fix it}

Remove the duplicate. A parameter given both positionally and by name also produces this.

\errorcode{E4171}{no parameter is named \errhole{}}
\label{sec:codes:e4171}

Raised during \textbf{validation}, from \texttt{ValidateErrorMessage::\allowbreak{}NO\_\allowbreak{}PARAMETER\_\allowbreak{}NAMED} (\textit{src/\allowbreak{}ymirc/\allowbreak{}semantic/\allowbreak{}validator/\allowbreak{}errors.yr}).

\subsubsection*{What it means}

A named argument at a call site names a parameter the callee does not have.

\subsubsection*{Example}

\textit{test\_\allowbreak{}resources/\allowbreak{}default\_\allowbreak{}ctor/\allowbreak{}test3.yr}:

%% check: skip
\begin{lstlisting}[style=coloredverbatimError]
// A field that is not public is not nameable at the call site, it is not part of
// the signature of the generated ctor.
class Foo {
   pub {
     let x: i32;
     let mut y: i32;
   }

   let mut z: i32 = 98;

   pub self;
}

fn main () {
    let a = copy Foo (y-> 2, 1, z-> 3);
    a.x;
}
\end{lstlisting}

\begin{lstlisting}[style=bashVerb, breaklines=true, escapechar=~]
Error[E4171] : no parameter is named z
 --> test_resources/default_ctor/test3.yr:(15,34)
15  ~\lstbox{\textSFxi{}}~     let a = copy Foo (y-> 2, 1, z-> 3);
    ~\lstbox{\textSFv{}}~                                  ^^
    ~\lstbox{\textSFii{}}~~\lstbox{\textSFx{}}~~\lstbox{\textSFx{}}~~\lstbox{\textSFx{}}~~\lstbox{\textSFx{}}~~\lstbox{\textSFx{}}~~\lstbox{\textSFx{}}~~\lstbox{\textSFx{}}~
\end{lstlisting}

\subsubsection*{How to fix it}

Check the spelling against the prototype. Note that a parameter with a default value can only be passed by name, so an unnamed parameter cannot be reached this way.

\errorcode{E4212}{\errhole{} parameters \errhole{} expected, but \errhole{} \errhole{} provided}
\label{sec:codes:e4212}

Raised during \textbf{validation}, from \texttt{ValidateErrorMessage::\allowbreak{}TOO\_\allowbreak{}FEW\_\allowbreak{}PARAMETERS} (\textit{src/\allowbreak{}ymirc/\allowbreak{}semantic/\allowbreak{}validator/\allowbreak{}errors.yr}).

\subsubsection*{What it means}

A call provides a different number of arguments than the callee takes. Both counts are in the message.

A parameter that declares a default value is filled by name or left to its default, so it is not part of the count a positional call is compared against -- \tokennolst{foo (x: i32, y: i32 = 8)} expects one argument, not two. When the extra arguments are explained by such a parameter, \hyperref[sec:codes:e4283]{E4283} hangs under this error and names it.

\subsubsection*{Example}

\textit{test\_\allowbreak{}resources/\allowbreak{}data\_\allowbreak{}record/\allowbreak{}ctors/\allowbreak{}test2.yr}:

%% check: skip
\begin{lstlisting}[style=coloredverbatimError]
@data
record Foo {
    let x: i32;
}

fn main() {
    let a = Foo();
    let b = Foo(x-> 34);
    a.x;
    b.x;
}
\end{lstlisting}

\begin{lstlisting}[style=bashVerb, breaklines=true, escapechar=~]
Error[E4212] : 1 parameters is expected, but 0 are provided
 --> test_resources/data_record/ctors/test2.yr:(7,16)
 7  ~\lstbox{\textSFxi{}}~     let a = Foo();
    ~\lstbox{\textSFv{}}~                ^
    ~\lstbox{\textSFii{}}~~\lstbox{\textSFx{}}~~\lstbox{\textSFx{}}~~\lstbox{\textSFx{}}~~\lstbox{\textSFx{}}~~\lstbox{\textSFx{}}~~\lstbox{\textSFx{}}~~\lstbox{\textSFx{}}~
\end{lstlisting}

\subsubsection*{How to fix it}

Pass the missing arguments. If the count looks right, check for a parameter with a default value: it is reachable by name only, so it is not one of the arguments a positional call is expected to provide.

\errorcode{E4226}{the call operator is not defined for \errhole{} and \errhole{}}
\label{sec:codes:e4226}

Raised during \textbf{validation}, from \texttt{ValidateErrorMessage::\allowbreak{}UNDEFINED\_\allowbreak{}CALL\_\allowbreak{}OP} (\textit{src/\allowbreak{}ymirc/\allowbreak{}semantic/\allowbreak{}validator/\allowbreak{}errors.yr}).

\subsubsection*{What it means}

A value was called with arguments no candidate accepts. The value and the arguments are in the message, and the candidates are listed under it.

\subsubsection*{Example}

\textit{test\_\allowbreak{}resources/\allowbreak{}diagnostics/\allowbreak{}test48.yr}:

%% check: skip
\begin{lstlisting}[style=coloredverbatimError]
fn main () {
    let x = 1;
    x ();
}
\end{lstlisting}

\begin{lstlisting}[style=bashVerb, breaklines=true, escapechar=~]
Error[E4226] : the call operator is not defined for x: i32 and {}
 --> test_resources/diagnostics/test48.yr:(3,7)
 3  ~\lstbox{\textSFxi{}}~     x ();
    ~\lstbox{\textSFv{}}~       ^
    ~\lstbox{\textSFxi{}}~ Error[E4107] : call operator is not defined for value 1
    ~\lstbox{\textSFxi{}}~  --> test_resources/diagnostics/test48.yr:(2,13)
    ~\lstbox{\textSFxi{}}~  2  ~\lstbox{\textSFxi{}}~     let x = 1;
    ~\lstbox{\textSFxi{}}~     ~\lstbox{\textSFv{}}~             ^
    ~\lstbox{\textSFxi{}}~     ~\lstbox{\textSFii{}}~~\lstbox{\textSFx{}}~~\lstbox{\textSFx{}}~~\lstbox{\textSFx{}}~~\lstbox{\textSFx{}}~~\lstbox{\textSFx{}}~~\lstbox{\textSFx{}}~~\lstbox{\textSFx{}}~
    ~\lstbox{\textSFii{}}~~\lstbox{\textSFx{}}~~\lstbox{\textSFx{}}~~\lstbox{\textSFx{}}~~\lstbox{\textSFx{}}~~\lstbox{\textSFx{}}~~\lstbox{\textSFvii{}}~~\lstbox{\textSFx{}}~
\end{lstlisting}

\subsubsection*{Notes}

A candidate is refused with a note under this diagnostic. Besides the notes of \hyperref[sec:codes:e4236]{E4236} for a template candidate, one carries no code of its own:

\begin{itemize}
\item \texttt{ref of type \textless{}\textgreater{} must be const, but is mutable} (formerly E4086): a mutable reference was passed to a \tokennolst{ref} parameter that is not \tokennolst{mut}.
\end{itemize}

\subsubsection*{How to fix it}

Match one of the listed candidates, cf. \hyperref[sec:codes:e4107]{E4107}.

\errorcode{E4246}{the creation of a delegate from a method has no effect}
\label{sec:codes:e4246}

Raised during \textbf{validation}, from \texttt{ValidateErrorMessage::\allowbreak{}UNECESSARY\_\allowbreak{}ADDRESS\_\allowbreak{}METHOD} (\textit{src/\allowbreak{}ymirc/\allowbreak{}semantic/\allowbreak{}validator/\allowbreak{}errors.yr}).

\subsubsection*{What it means}

A delegate was taken of a method in a place where the delegate is used immediately, so building it changes nothing.

\subsubsection*{Example}

\textit{test\_\allowbreak{}resources/\allowbreak{}diagnostics/\allowbreak{}test78.yr}:

%% check: skip
\begin{lstlisting}[style=coloredverbatimError]
class Foo {
    pub self () {}
    pub fn bar (self) {}
}

fn main () {
    let f = copy Foo ();
    (&f.bar) ();
}
\end{lstlisting}

\begin{lstlisting}[style=bashVerb, breaklines=true, escapechar=~]
Error[E4246] : the creation of a delegate from a method has no effect
 --> test_resources/diagnostics/test78.yr:(8,6)
 8  ~\lstbox{\textSFxi{}}~     (&f.bar) ();
    ~\lstbox{\textSFv{}}~      ^
    ~\lstbox{\textSFii{}}~~\lstbox{\textSFx{}}~~\lstbox{\textSFx{}}~~\lstbox{\textSFx{}}~~\lstbox{\textSFx{}}~~\lstbox{\textSFx{}}~~\lstbox{\textSFx{}}~~\lstbox{\textSFx{}}~
\end{lstlisting}

\subsubsection*{How to fix it}

Call the method directly.

\errorcode{E4250}{copy the lambda closure has no effect, no values are enclosed}
\label{sec:codes:e4250}

Raised during \textbf{validation}, from \texttt{ValidateErrorMessage::\allowbreak{}UNECESSARY\_\allowbreak{}LMBD\_\allowbreak{}COPY} (\textit{src/\allowbreak{}ymirc/\allowbreak{}semantic/\allowbreak{}validator/\allowbreak{}errors.yr}).

\subsubsection*{What it means}

\tokennolst{copy} was applied to a lambda that captures nothing, so there is no environment to duplicate.

\subsubsection*{Example}

\textit{test\_\allowbreak{}resources/\allowbreak{}diagnostics/\allowbreak{}test79.yr}:

%% check: skip
\begin{lstlisting}[style=coloredverbatimError]
fn main () {
    let f = copy (|| => { 1 });
    f ();
}
\end{lstlisting}

\begin{lstlisting}[style=bashVerb, breaklines=true, escapechar=~]
Error[E4250] : copy the lambda closure has no effect, no values are enclosed
 --> test_resources/diagnostics/test79.yr:(2,13)
 2  ~\lstbox{\textSFxi{}}~     let f = copy (|| => { 1 });
    ~\lstbox{\textSFv{}}~             ^^^^
    ~\lstbox{\textSFii{}}~~\lstbox{\textSFx{}}~~\lstbox{\textSFx{}}~~\lstbox{\textSFx{}}~~\lstbox{\textSFx{}}~~\lstbox{\textSFx{}}~~\lstbox{\textSFx{}}~~\lstbox{\textSFx{}}~
\end{lstlisting}

\subsubsection*{How to fix it}

Drop the \tokennolst{copy}.

\errorcode{E4280}{\errhole{} called with \errhole{} works with multiple candidates}
\label{sec:codes:e4280}

Raised during \textbf{validation}, from \texttt{ValidateErrorMessage::\allowbreak{}WORKS\_\allowbreak{}WITH\_\allowbreak{}BOTH} (\textit{src/\allowbreak{}ymirc/\allowbreak{}semantic/\allowbreak{}validator/\allowbreak{}errors.yr}).

\subsubsection*{What it means}

A call matches several candidates equally well, so overload resolution has no single answer. The candidates are listed under the diagnostic.

\subsubsection*{Example}

\textit{test\_\allowbreak{}resources/\allowbreak{}root\_\allowbreak{}path/\allowbreak{}test7.yr}:

%% check: skip
\begin{lstlisting}[style=coloredverbatimError]
mod A {
    pub fn foo ()-> i32 { 12 }
}

mod B {
    mod A {
        pub fn foo ()-> f64 { 12.0 }
    }

    pub fn bar ()-> f64 {
        A::foo ()
    }
}

fn main () {
    let a = B::bar ();
    a;
}
\end{lstlisting}

\begin{lstlisting}[style=bashVerb, breaklines=true, escapechar=~]
Error[E4280] : foo called with {} works with multiple candidates
 --> test_resources/root_path/test7.yr:(11,16)
11  ~\lstbox{\textSFxi{}}~         A::foo ()
    ~\lstbox{\textSFv{}}~                ^  note: candidate test7::A::foo ()-> i32 -- (test_resources/root_path/test7.yr:(2,12))
    ~\lstbox{\textSFxi{}}~                ~\lstbox{\textSFii{}}~~\lstbox{\textSFx{}}~ note: candidate test7::B::A::foo ()-> f64 -- (test_resources/root_path/test7.yr:(7,16))
    ~\lstbox{\textSFii{}}~~\lstbox{\textSFx{}}~~\lstbox{\textSFx{}}~~\lstbox{\textSFx{}}~~\lstbox{\textSFx{}}~~\lstbox{\textSFx{}}~~\lstbox{\textSFx{}}~~\lstbox{\textSFx{}}~
\end{lstlisting}

\subsubsection*{How to fix it}

Make the call select one: convert an argument explicitly, pass an argument by name, or qualify the symbol with its module.

\errorcode{E4283}{the parameter \errhole{} has a default value, it can only be passed by name}
\label{sec:codes:e4283}

Raised during \textbf{validation}, from \texttt{ValidateErrorMessage::\allowbreak{}DEFAULT\_\allowbreak{}PARAM\_\allowbreak{}BY\_\allowbreak{}POSITION} (\textit{src/\allowbreak{}ymirc/\allowbreak{}semantic/\allowbreak{}validator/\allowbreak{}errors.yr}).

\subsubsection*{What it means}

A parameter that declares a default value is filled by name or left to its default -- it is never a slot a positional argument lands in. A call that provides one argument too many for such a prototype was written expecting the opposite, so this hangs under the \hyperref[sec:codes:e4212]{E4212} arity error and names the parameter responsible for the difference between the two counts.

It is not raised when a call is \textit{short} of an argument: the missing parameter is then one that has no default, and naming a defaulted one would point at the wrong place.

A tuple cast to a record is expanded into a positional call, so it reaches this error too. A tuple has no name syntax, so the fix there is to drop the element and let the field take its default.

\subsubsection*{Example}

\textit{test\_\allowbreak{}resources/\allowbreak{}diagnostics/\allowbreak{}test83.yr}:

%% check: skip
\begin{lstlisting}[style=coloredverbatimError]
fn foo (x : i32, y : i32 = 8)-> i32 {
    x + y
}

fn main () {
    foo (1, 2);
}
\end{lstlisting}

\begin{lstlisting}[style=bashVerb, breaklines=true, escapechar=~]
Error[E4212] : 1 parameters is expected, but 2 are provided
 --> test_resources/diagnostics/test83.yr:(6,9)
 6  ~\lstbox{\textSFxi{}}~     foo (1, 2);
    ~\lstbox{\textSFv{}}~         ^
    ~\lstbox{\textSFxi{}}~ Error[E4283] : the parameter y has a default value, it can only be passed by name
    ~\lstbox{\textSFxi{}}~  --> test_resources/diagnostics/test83.yr:(1,18)
    ~\lstbox{\textSFxi{}}~  1  ~\lstbox{\textSFxi{}}~ fn foo (x : i32, y : i32 = 8)-> i32 {
    ~\lstbox{\textSFxi{}}~     ~\lstbox{\textSFv{}}~                  ^
    ~\lstbox{\textSFxi{}}~     ~\lstbox{\textSFii{}}~~\lstbox{\textSFx{}}~~\lstbox{\textSFx{}}~~\lstbox{\textSFx{}}~~\lstbox{\textSFx{}}~~\lstbox{\textSFx{}}~~\lstbox{\textSFx{}}~~\lstbox{\textSFx{}}~
    ~\lstbox{\textSFii{}}~~\lstbox{\textSFx{}}~~\lstbox{\textSFx{}}~~\lstbox{\textSFx{}}~~\lstbox{\textSFx{}}~~\lstbox{\textSFx{}}~~\lstbox{\textSFvii{}}~~\lstbox{\textSFx{}}~
\end{lstlisting}

\subsubsection*{How to fix it}

Pass the parameter by name:

%% check: skip
\begin{lstlisting}[style=coloredverbatim]
foo (1, y-> 2);
\end{lstlisting}

or drop the argument and let the default apply:

%% check: skip
\begin{lstlisting}[style=coloredverbatim]
foo (1);
\end{lstlisting}

\errorcode{E4284}{the parameter \errhole{} can only be passed by name}
\label{sec:codes:e4284}

Raised during \textbf{validation}, from \texttt{ValidateErrorMessage::\allowbreak{}NAMED\_\allowbreak{}PARAM\_\allowbreak{}BY\_\allowbreak{}POSITION} (\textit{src/\allowbreak{}ymirc/\allowbreak{}semantic/\allowbreak{}validator/\allowbreak{}errors.yr}).

\subsubsection*{What it means}

A parameter marked \tokennolst{@} is forced named: a call fills it by name, in any position relative to the positional arguments, and never by position. A call that leaves it unnamed is rejected, one error per such parameter, each pointing at the parameter's declaration.

When the positional arguments are otherwise the right number, this is the only reason the call fails. When they are not, it comes with the \hyperref[sec:codes:e4212]{E4212} arity error it explains: the positional values that were written for the marked parameters are the ones counted too many or too few.

The call is rejected as a candidate, so these errors are listed under the candidate note of \hyperref[sec:codes:e4226]{E4226}.

\subsubsection*{Example}

\textit{test\_\allowbreak{}resources/\allowbreak{}diagnostics/\allowbreak{}test145.yr}:

%% check: skip
\begin{lstlisting}[style=coloredverbatimError]
fn foo (@a: i32, b: i32)-> i32 { a + b }

fn main () { foo (1); }
\end{lstlisting}

\begin{lstlisting}[style=bashVerb, breaklines=true, escapechar=~]
Error[E4226] : the call operator is not defined for test145::foo (@a: i32, b: i32)-> i32 and {1: i32}
 --> test_resources/diagnostics/test145.yr:(3,18)
 1  ~\lstbox{\textSFxi{}}~ fn foo (@a: i32, b: i32)-> i32 { a + b }
    ~\lstbox{\textSFv{}}~    ---  note: candidate test145::foo (@a: i32, b: i32)-> i32
    ~\lstbox{\textSFxi{}}~      ~\lstbox{\textSFii{}}~~\lstbox{\textSFx{}}~ error: the parameter a can only be passed by name
 2  ~\lstbox{\textSFxi{}}~
 3  ~\lstbox{\textSFxi{}}~ fn main () { foo (1); }
    ~\lstbox{\textSFv{}}~                  ^
    ~\lstbox{\textSFxi{}}~
    = when validating test145::main -- (test_resources/diagnostics/test145.yr:(3,4))
    ~\lstbox{\textSFii{}}~~\lstbox{\textSFx{}}~~\lstbox{\textSFx{}}~~\lstbox{\textSFx{}}~~\lstbox{\textSFx{}}~~\lstbox{\textSFx{}}~~\lstbox{\textSFx{}}~~\lstbox{\textSFx{}}~
\end{lstlisting}

\subsubsection*{How to fix it}

Name the parameter at the call site:

%% check: skip
\begin{lstlisting}[style=coloredverbatim]
foo (a-> 1, 2);
\end{lstlisting}

\errorcode{E4285}{parameter with a default value cannot be marked \errhole{}, it can already only be passed by name}
\label{sec:codes:e4285}

Raised during \textbf{validation}, from \texttt{ValidateErrorMessage::\allowbreak{}NAMED\_\allowbreak{}PARAM\_\allowbreak{}WITH\_\allowbreak{}DEFAULT} (\textit{src/\allowbreak{}ymirc/\allowbreak{}semantic/\allowbreak{}validator/\allowbreak{}errors.yr}).

\subsubsection*{What it means}

The \tokennolst{@} marker forces a parameter to be passed by name. A parameter that declares a default value is already filled by name or left to its default, never by position (cf. \hyperref[sec:codes:e4283]{E4283}), so the marker has nothing left to ask and the two cannot be combined. The error points at the parameter and at its default value.

\subsubsection*{Example}

\textit{test\_\allowbreak{}resources/\allowbreak{}diagnostics/\allowbreak{}test84.yr}:

%% check: skip
\begin{lstlisting}[style=coloredverbatimError]
fn foo (@a: i32 = 2)-> i32 { a }

fn main () { foo (); }
\end{lstlisting}

\begin{lstlisting}[style=bashVerb, breaklines=true, escapechar=~]
Error[E4285] : parameter with a default value cannot be marked @, it can already only be passed by name
 --> test_resources/diagnostics/test84.yr:(1,10)
 1  ~\lstbox{\textSFxi{}}~ fn foo (@a: i32 = 2)-> i32 { a }
    ~\lstbox{\textSFv{}}~          ^        -
    ~\lstbox{\textSFxi{}}~
    = when validating test84::foo -- (test_resources/diagnostics/test84.yr:(1,4))
    ~\lstbox{\textSFii{}}~~\lstbox{\textSFx{}}~~\lstbox{\textSFx{}}~~\lstbox{\textSFx{}}~~\lstbox{\textSFx{}}~~\lstbox{\textSFx{}}~~\lstbox{\textSFx{}}~~\lstbox{\textSFx{}}~
Error[E4104] : symbol test84::foo has errors
 --> test_resources/diagnostics/test84.yr:(3,14)
 1  ~\lstbox{\textSFxi{}}~ fn foo (@a: i32 = 2)-> i32 { a }
    ~\lstbox{\textSFv{}}~          -        -  error: parameter with a default value cannot be marked @, it can already only be passed by name
 2  ~\lstbox{\textSFxi{}}~
 3  ~\lstbox{\textSFxi{}}~ fn main () { foo (); }
    ~\lstbox{\textSFv{}}~              ^^^
    ~\lstbox{\textSFxi{}}~
    = when validating test84::foo -- (test_resources/diagnostics/test84.yr:(1,4))
    = when validating test84::main -- (test_resources/diagnostics/test84.yr:(3,4))
    ~\lstbox{\textSFii{}}~~\lstbox{\textSFx{}}~~\lstbox{\textSFx{}}~~\lstbox{\textSFx{}}~~\lstbox{\textSFx{}}~~\lstbox{\textSFx{}}~~\lstbox{\textSFx{}}~~\lstbox{\textSFx{}}~
\end{lstlisting}

\subsubsection*{How to fix it}

Drop the marker, the parameter stays passed by name:

%% check: skip
\begin{lstlisting}[style=coloredverbatim]
fn foo (a: i32 = 2)-> i32 { a }
\end{lstlisting}

or drop the default value, if the parameter has to be provided at every call.

\errorcode{E4298}{the method \errhole{} is not a value, it can only be called}
\label{sec:codes:e4298}

Raised during \textbf{validation}, from \texttt{ValidateErrorMessage::\allowbreak{}FAKE\_\allowbreak{}METHOD\_\allowbreak{}CALL} (\textit{src/\allowbreak{}ymirc/\allowbreak{}semantic/\allowbreak{}validator/\allowbreak{}errors.yr}).

\subsubsection*{What it means}

A fake method is not a field, it is a method syntax the compiler rewrites when it is called. There is nothing to refer to before that rewriting, so the access cannot be used as a value of its own. They are \tokennolst{map:.remove}, \tokennolst{range.stepBy}, \tokennolst{range.reverse} and \tokennolst{generator.next}.

\subsubsection*{Example}

\textit{test\_\allowbreak{}resources/\allowbreak{}generators/\allowbreak{}test32.yr}:

%% check: skip
\begin{lstlisting}[style=coloredverbatimError]
// the fake method next is not a value
yield range (i: i32, j: i32)-> i32 {
    for n in i .. j {
        yield n;
    }
}

fn main () {
    let r = range (0, 2);
    let n = r.next;
    n;
}
\end{lstlisting}

\begin{lstlisting}[style=bashVerb, breaklines=true, escapechar=~]
Error[E4298] : the method next is not a value, it can only be called
 --> test_resources/generators/test32.yr:(10,14)
10  ~\lstbox{\textSFxi{}}~     let n = r.next;
    ~\lstbox{\textSFv{}}~              ^----
    ~\lstbox{\textSFxi{}}~
    = when validating test32::main -- (test_resources/generators/test32.yr:(8,4))
    ~\lstbox{\textSFii{}}~~\lstbox{\textSFx{}}~~\lstbox{\textSFx{}}~~\lstbox{\textSFx{}}~~\lstbox{\textSFx{}}~~\lstbox{\textSFx{}}~~\lstbox{\textSFx{}}~~\lstbox{\textSFx{}}~
\end{lstlisting}

\subsubsection*{How to fix it}

Call the method. \tokennolst{r.next ()} returns the next value of the generator wrapped in an option, \tokennolst{none} once the generator is over. \tokennolst{r:.next ()} borrows that value mutably instead of copying it constantly, on a mutable generator yielding borrowed data.
